\chapter{System Implementation}

The project is implemented as a three-component privacy-preserving authentication system. The implementation involves cryptographic circuit design, mobile app development, server-side proof verification, and web SDK integration. This chapter describes each component's implementation in detail with code listings and architectural insights.

\section{Development Environment Setup}

Before beginning the core implementation, a consistent development environment was established across all components:

\begin{lstlisting}[language=bash, caption={Development Environment Initialization}]
# Install Node.js 18 LTS
nvm install 18
nvm use 18

# Install Expo CLI for mobile development
npm install -g expo-cli eas-cli

# Install Circom compiler
cargo install circom

# Install snarkjs globally
npm install -g snarkjs

# Verify installations
node --version   # v18.x.x
circom --version # circom compiler 2.0.x
snarkjs --version  # 0.7.x
\end{lstlisting}

\section{Wallet Application}

The Wallet is a React Native/Expo mobile application that serves as the user's identity container and proof generator. Its file structure is organized as follows:

\begin{lstlisting}[language=bash, caption={Wallet Application File Structure}]
wallet/
  app/
    (tabs)/
      index.tsx       # Credentials home screen
      sessions.tsx    # Session history screen
      settings.tsx    # Settings screen
  components/
    CredentialCard.tsx
    ProofModal.tsx
    QRScanner.tsx
  hooks/
    useCredentialStore.ts
    useNetworkStatus.ts
    useProofEngine.ts
  lib/
    wallet.ts         # Ed25519 identity
    zkbridge.ts       # WebView ZK proof bridge
    crypto.ts         # Hashing utilities
    revocation.ts     # Revocation checking
  circuits/
    age_check.wasm
    age_check.zkey
    student_check.wasm
    student_check.zkey
  assets/
    zkwebview.html    # Sandboxed ZK environment
\end{lstlisting}

\subsection{Identity Management (Ed25519)}

The wallet generates an Ed25519 keypair on first launch for decentralized identity:

\begin{lstlisting}[language=Java, caption={Identity Generation (wallet.ts)}]
export async function generateAndStoreIdentity() {
    const privateKey = new Uint8Array(32);
    getRandomValues(privateKey);
    const publicKey = ed25519.getPublicKey(privateKey);
    
    // Store securely in platform SecureStore
    await SecureStore.setItemAsync(
        PRIVATE_KEY_ALIAS, toBase64(privateKey)
    );
    
    // Derive W3C did:key identifier
    const did = deriveDID(publicKey); // "did:key:z..."
    return { publicKey, did };
}

function deriveDID(publicKey: Uint8Array): string {
    // Prepend multicodec prefix 0xed01
    const multicodec = new Uint8Array([0xed, 0x01, ...publicKey]);
    // Encode with Base58BTC
    const b58 = base58btc.encode(multicodec);
    return `did:key:z${b58}`;
}
\end{lstlisting}

The DID is derived following the \textbf{W3C did:key specification}: the Ed25519 public key is prepended with the multicodec prefix \texttt{0xed01}, encoded in Base58BTC, and prefixed with \texttt{did:key:z}.

\subsection{Credential Storage}

Credentials are stored locally using Zustand state management with AsyncStorage persistence. Table 5.1 below describes each field in the credential data structure stored within the wallet.

\begin{table}[H]
\centering
\caption{Credential Structure}
\begin{tabular}{|l|p{8cm}|}
\hline
\textbf{Field} & \textbf{Description} \\
\hline
\texttt{id} & Unique identifier (cryptographically generated) \\
\texttt{type} & Credential type (Age Verification, Student ID, Trial) \\
\texttt{issuer} & Name of the issuing authority \\
\texttt{attributes} & Key-value pairs (e.g., birth\_year: ``1995'') \\
\texttt{commitments} & Poseidon hash commitments binding attributes \\
\texttt{expiresAt} & Optional expiration timestamp \\
\texttt{revocationId} & Optional ID for revocation checking \\
\hline
\end{tabular}
\end{table}

\begin{lstlisting}[language=Java, caption={Credential Store (useCredentialStore.ts)}]
interface Credential {
  id: string;
  type: 'Age Verification' | 'Student ID' | 'Trial';
  issuer: string;
  attributes: Record<string, string | number | boolean>;
  commitments: Record<string, bigint>;
  expiresAt?: number;
  revocationId?: string;
}

const useCredentialStore = create<CredentialStore>()(
  persist(
    (set, get) => ({
      credentials: [],
      addCredential: (cred: Credential) =>
        set(s => ({ credentials: [...s.credentials, cred] })),
      removeCredential: (id: string) =>
        set(s => ({
          credentials: s.credentials.filter(c => c.id !== id)
        })),
    }),
    { name: 'credentials', storage: createJSONStorage(() => AsyncStorage) }
  )
);
\end{lstlisting}

\subsection{PIN-based Protection}

The wallet is protected by a PIN with security-hardened storage using SHA-256 hashing with unique salts. PIN verification uses \textbf{timing-safe comparison} to prevent timing side-channel attacks:

\begin{lstlisting}[language=Java, caption={Timing-safe PIN verification (utils.ts)}]
export async function verifyPin(
    pin: string, storedValue: string
): Promise<boolean> {
    const [salt, storedHash] = storedValue.split(':');
    const inputHash = await digestStringAsync(
        CryptoDigestAlgorithm.SHA256, pin + salt
    );
    // Timing-safe comparison
    let result = 0;
    for (let i = 0; i < inputHash.length; i++) {
        result |= inputHash.charCodeAt(i) 
                 ^ storedHash.charCodeAt(i);
    }
    return result === 0;
}

export async function hashPin(pin: string): Promise<string> {
    // Generate random 16-byte salt
    const salt = toHex(getRandomValues(new Uint8Array(16)));
    const hash = await digestStringAsync(
        CryptoDigestAlgorithm.SHA256, pin + salt
    );
    return `${salt}:${hash}`;
}
\end{lstlisting}

\subsection{QR Scanner and Session Parsing}

When a user opens the Wallet and taps ``Scan QR Code'', the camera module activates and scans the QR code displayed on the verifier website:

\begin{lstlisting}[language=Java, caption={QR Code Parsing (QRScanner.tsx)}]
function parseQRPayload(rawData: string): QRPayload {
  const decoded = atob(rawData);
  const payload: QRPayload = JSON.parse(decoded);
  
  // Validate protocol version
  if (payload.v !== 1) throw new Error('Unsupported protocol version');
  
  // Check expiration
  const now = Math.floor(Date.now() / 1000);
  if (payload.expires_at < now) throw new Error('QR code expired');
  
  // Validate required fields
  if (!payload.session_id || !payload.nonce || 
      !payload.credential_type) {
    throw new Error('Invalid QR payload structure');
  }
  return payload;
}
\end{lstlisting}

\section{ZK Circuit Design}

\subsection{Age Verification Circuit}

The \texttt{age\_check.circom} circuit implements age verification with Poseidon commitment binding:

\begin{lstlisting}[language=C, caption={Age Verification Circuit (age\_check.circom)}]
pragma circom 2.0.0;

include "node_modules/circomlib/circuits/poseidon.circom";
include "node_modules/circomlib/circuits/comparators.circom";

template AgeCheck() {
    signal input currentYear;    // Public
    signal input minAge;         // Public
    signal input birthYear;      // Private
    signal input salt;           // Private
    signal input commitment;     // Public

    // 1. Verify Commitment: commitment = Poseidon(birthYear, salt)
    component hasher = Poseidon(2);
    hasher.inputs[0] <== birthYear;
    hasher.inputs[1] <== salt;
    hasher.out === commitment;

    // 2. Age Check: currentYear - birthYear >= minAge
    signal age;
    age <== currentYear - birthYear;
    component ge = GreaterEqThan(32);
    ge.in[0] <== age;
    ge.in[1] <== minAge;
    ge.out === 1;
}

component main {public [currentYear, minAge, commitment]} = AgeCheck();
\end{lstlisting}

\subsection{Student ID Verification Circuit}

The \texttt{student\_check.circom} circuit verifies active student status using a 3-input Poseidon hash commitment:

\begin{lstlisting}[language=C, caption={Student ID Circuit (student\_check.circom)}]
pragma circom 2.0.0;

include "node_modules/circomlib/circuits/poseidon.circom";
include "node_modules/circomlib/circuits/comparators.circom";

template StudentCheck() {
    signal input currentYear;    // Public
    signal input commitment;     // Public
    signal input isStudent;      // Private (must be 1)
    signal input expiryYear;     // Private
    signal input salt;           // Private

    // 1. Verify commitment
    component hasher = Poseidon(3);
    hasher.inputs[0] <== isStudent;
    hasher.inputs[1] <== expiryYear;
    hasher.inputs[2] <== salt;
    hasher.out === commitment;

    // 2. isStudent === 1
    isStudent === 1;

    // 3. expiryYear >= currentYear (not expired)
    component ge = GreaterEqThan(32);
    ge.in[0] <== expiryYear;
    ge.in[1] <== currentYear;
    ge.out === 1;
}

component main {public [currentYear, commitment]} = StudentCheck();
\end{lstlisting}

\subsection{Circuit Compilation and Trusted Setup}

\begin{lstlisting}[language=bash, caption={Circuit Compilation Steps}]
# Compile circuit to R1CS and WASM
circom age_check.circom --r1cs --wasm --sym

# Download Powers of Tau (Phase 1 trusted setup)
wget https://hermez.s3-eu-west-1.amazonaws.com/powersOfTau28_hez_final_14.ptau

# Phase 2 setup for circuit
snarkjs groth16 setup age_check.r1cs \
    powersOfTau28_hez_final_14.ptau \
    age_check_0.zkey

# Contribute to trusted setup (add randomness)
snarkjs zkey contribute age_check_0.zkey \
    age_check_final.zkey --name="ZeroAuth Setup"

# Export verification key
snarkjs zkey export verificationkey \
    age_check_final.zkey verification_key.json

# Export Solidity verifier (for future on-chain use)
snarkjs zkey export solidityverifier \
    age_check_final.zkey AgeVerifier.sol
\end{lstlisting}

\subsection{WebView ZK Proof Bridge}

Since snarkjs runs in a browser environment (WebView), a JavaScript bridge connects the React Native app to the proof generation engine:

\begin{lstlisting}[language=Java, caption={ZK WebView Bridge (zkbridge.ts)}]
export async function generateProof(
    credType: string,
    inputs: ProofInputs,
    timeout = 90000
): Promise<ZKProofPayload> {
    return new Promise((resolve, reject) => {
        const timer = setTimeout(
            () => reject(new Error('Proof timeout')), timeout
        );
        
        webViewRef.current?.injectJavaScript(`
            (async () => {
                try {
                    const { proof, publicSignals } = 
                        await snarkjs.groth16.fullProve(
                            ${JSON.stringify(inputs)},
                            '${credType}_check.wasm',
                            '${credType}_check.zkey'
                        );
                    window.ReactNativeWebView.postMessage(
                        JSON.stringify({ success: true, 
                            proof, publicSignals })
                    );
                } catch (e) {
                    window.ReactNativeWebView.postMessage(
                        JSON.stringify({ success: false, 
                            error: e.message })
                    );
                }
            })();
        `);
        
        // onMessage handler resolves/rejects the promise
        setOnMessage((data) => {
            clearTimeout(timer);
            if (data.success) resolve(data);
            else reject(new Error(data.error));
        });
    });
}
\end{lstlisting}

\section{Relay Server}

The Relay is a Node.js/Express server that manages verification sessions and verifies ZK proofs.

\subsection{Server Architecture}

\begin{lstlisting}[language=bash, caption={Relay Server File Structure}]
relay/
  src/
    index.ts           # Express app entry point
    routes/
      sessions.ts      # Session CRUD + proof submission
    middleware/
      rateLimiter.ts   # express-rate-limit config
      errorHandler.ts  # Global error middleware
      logger.ts        # Request logging
    services/
      supabase.ts      # Database client
      zkVerifier.ts    # snarkjs proof verification
      sessionService.ts # Session lifecycle management
    utils/
      proofHash.ts     # SHA-256 proof hashing
      cleanup.ts       # Expired session cleanup job
\end{lstlisting}

\subsection{Session Management}

Sessions follow a lifecycle: \texttt{PENDING} $\rightarrow$ \texttt{COMPLETED} $\rightarrow$ \texttt{expired/deleted}. New sessions are created with a 5-minute TTL. Expired sessions are automatically cleaned up every 60 seconds.

\begin{lstlisting}[language=Java, caption={Session Creation (sessionService.ts)}]
export async function createSession(
    credentialType: string,
    verifierName: string,
    requiredClaims: string[],
    expirySeconds = 300
): Promise<Session> {
    const sessionId = uuidv4();
    const nonce = crypto.randomBytes(32).toString('hex');
    const expiresAt = new Date(
        Date.now() + expirySeconds * 1000
    ).toISOString();
    
    const { data, error } = await supabase
        .from('sessions')
        .insert({
            session_id: sessionId,
            nonce,
            verifier_name: verifierName,
            required_claims: requiredClaims,
            credential_type: credentialType,
            status: 'PENDING',
            expires_at: expiresAt
        })
        .select()
        .single();
    
    if (error) throw new Error(`DB insert failed: ${error.message}`);
    return data;
}
\end{lstlisting}

\subsection{Proof Verification Pipeline}

When a proof is submitted, it passes through a multi-stage verification pipeline:

\begin{enumerate}
    \item \textbf{Session Validation}: Check session exists and is PENDING.
    \item \textbf{Schema Validation}: Verify proof structure ($\pi_a$, $\pi_b$, $\pi_c$ arrays).
    \item \textbf{Replay Protection}: Compute SHA-256 hash and check for duplicates.
    \item \textbf{ZK Verification}: Cryptographically verify using \texttt{snarkjs.groth16.verify()}.
    \item \textbf{Session Update}: Mark session as COMPLETED with proof data.
\end{enumerate}

\begin{lstlisting}[language=Java, caption={Proof Verification (zkVerifier.ts)}]
export async function verifyProof(
    sessionId: string,
    proof: ZKProof,
    publicSignals: string[]
): Promise<boolean> {
    // 1. Load verification key from database
    const { data: keyData } = await supabase
        .from('verification_keys')
        .select('key_data')
        .eq('credential_type', session.credential_type)
        .single();
    
    // 2. Compute proof hash for replay check
    const proofHash = computeProofHash(proof, publicSignals);
    if (await isProofUsed(proofHash)) {
        throw new Error('Proof replay detected');
    }
    
    // 3. Verify with snarkjs
    const isValid = await snarkjs.groth16.verify(
        keyData.key_data,
        publicSignals,
        proof
    );
    
    if (!isValid) throw new Error('Invalid ZK proof');
    
    // 4. Mark session completed
    await supabase.from('sessions')
        .update({ 
            status: 'COMPLETED', 
            proof, proof_hash: proofHash 
        })
        .eq('session_id', sessionId);
    
    return true;
}

function computeProofHash(
    proof: ZKProof, 
    publicSignals: string[]
): string {
    const canonical = JSON.stringify({
        pi_a: proof.pi_a.sort(),
        pi_b: proof.pi_b.map(p => p.sort()).sort(),
        pi_c: proof.pi_c.sort(),
        publicSignals: publicSignals.sort()
    });
    return crypto.createHash('sha256')
        .update(canonical).digest('hex');
}
\end{lstlisting}

\subsection{Security Measures}

The Relay server incorporates multiple independent security mechanisms to protect against common attack vectors. Table 5.2 below summarises each measure and its corresponding implementation.

\begin{table}[H]
\centering
\caption{Relay Server Security Measures}
\begin{tabular}{|l|p{7cm}|}
\hline
\textbf{Security Measure} & \textbf{Implementation} \\
\hline
Rate Limiting & 500 requests per IP per 15 minutes \\
Request Size Limit & 10MB max body, 1MB max proof payload \\
Session Expiry & 5-minute TTL with automatic cleanup \\
Proof Replay Protection & SHA-256 hash deduplication \\
Row-Level Security & Supabase RLS policies on all tables \\
CORS & Enabled for cross-origin SDK access \\
Request Logging & Timestamp, method, path, IP logged \\
Error Sanitization & Internal errors never exposed to client \\
HTTPS Only & HTTP redirected to HTTPS in production \\
\hline
\end{tabular}
\end{table}

\section{Web SDK}

The SDK provides two integration methods for verifier websites:

\subsection{React Component Integration}
\begin{lstlisting}[language=HTML, caption={React Integration}]
import { ZeroAuthButton } from '@zero-auth/sdk';

// Basic usage
<ZeroAuthButton 
  credentialType="Age Verification"
  claims={['birth_year']}
  onSuccess={(result) => {
    console.log('Verified!', result);
    // Grant access to user
    setVerified(true);
  }}
  onError={(err) => {
    console.error('Verification failed:', err);
  }}
/>

// With custom relay URL
<ZeroAuthButton 
  relayUrl="https://your-relay.example.com"
  credentialType="Student ID"
  claims={['isStudent', 'expiryYear']}
  onSuccess={handleSuccess}
  verifierName="University Bookstore"
/>
\end{lstlisting}

\subsection{CDN / Vanilla JavaScript Integration}
\begin{lstlisting}[language=HTML, caption={CDN Integration}]
<script src="https://unpkg.com/@zero-auth/sdk/dist/
  zero-auth-sdk.umd.js"></script>
<script>
  window.ZeroAuthConfig = {
    relayUrl: 'https://zeroauth-relay.onrender.com',
    verifierName: 'My Website'
  };
  
  ZeroAuth.createButton({
    elementId: 'verify-btn',
    credentialType: 'Age Verification',
    claims: ['birth_year'],
    onSuccess: function(result) {
      document.getElementById('content').style.display = 'block';
    }
  });
</script>

<div id="verify-btn"></div>
<div id="content" style="display:none">
  <!-- Age-restricted content -->
</div>
\end{lstlisting}

\subsection{SDK Session Polling Implementation}

\begin{lstlisting}[language=Java, caption={Session Polling (sdk/polling.ts)}]
export async function pollSessionStatus(
    sessionId: string,
    relayUrl: string,
    intervalMs = 2000,
    timeoutMs = 300000
): Promise<SessionResult> {
    const startTime = Date.now();
    
    while (Date.now() - startTime < timeoutMs) {
        const res = await fetch(
            `${relayUrl}/api/v1/sessions/${sessionId}`
        );
        const data = await res.json();
        
        if (data.status === 'COMPLETED') {
            return { success: true, ...data };
        }
        if (data.status === 'EXPIRED' || data.status === 'REJECTED') {
            throw new Error(`Session ${data.status}`);
        }
        
        // Wait before next poll
        await new Promise(r => setTimeout(r, intervalMs));
    }
    throw new Error('Verification timeout');
}
\end{lstlisting}

\section{QR Code Protocol}

The verification protocol uses a structured JSON payload encoded in QR codes:

\begin{lstlisting}[language=Java, caption={QR Code Payload Structure}]
{
  "v": 1,
  "action": "verify",
  "session_id": "550e8400-e29b-41d4-a716-446655440000",
  "nonce": "a3f2...<64 hex chars>",
  "verifier": {
    "name": "Example Age-Gated Service",
    "did": "did:web:example.com",
    "callback": "https://relay/api/v1/sessions/{id}/proof"
  },
  "required_claims": ["birth_year"],
  "credential_type": "Age Verification",
  "expires_at": 1712345678
}
\end{lstlisting}

\section{Security Implementation Details}

\subsection{Proof Replay Protection}

Each proof is hashed using SHA-256 with deterministic key ordering. If a proof hash matches a previously submitted proof for the same session, the submission is rejected.

\subsection{Credential Revocation}

The wallet supports offline-capable revocation checking with a 1-hour cache duration. Revocation status is checked before generating proofs and cached in AsyncStorage for offline scenarios.

\subsection{Offline Queue}

When network connectivity is unavailable, actions are queued and automatically processed upon reconnection using the \texttt{useNetworkStatus} hook that monitors connectivity changes via \texttt{NetInfo}.

\begin{lstlisting}[language=Java, caption={Offline Queue (useNetworkStatus.ts)}]
export function useNetworkStatus() {
  const [isOnline, setIsOnline] = useState(true);
  const queueRef = useRef<QueuedAction[]>([]);

  useEffect(() => {
    const unsub = NetInfo.addEventListener(state => {
      const online = state.isConnected && state.isInternetReachable;
      if (online && !isOnline) {
        // Process queued actions on reconnect
        processQueue(queueRef.current).then(() => {
          queueRef.current = [];
        });
      }
      setIsOnline(online ?? false);
    });
    return unsub;
  }, [isOnline]);

  const enqueue = (action: QueuedAction) => {
    queueRef.current.push(action);
  };

  return { isOnline, enqueue };
}
\end{lstlisting}

\section{Database Implementation}

The Supabase PostgreSQL schema is set up with the following SQL statements:

\begin{lstlisting}[language=SQL, caption={Database Schema DDL}]
-- Sessions table
CREATE TABLE sessions (
    id UUID PRIMARY KEY DEFAULT gen_random_uuid(),
    session_id TEXT UNIQUE NOT NULL,
    nonce TEXT NOT NULL,
    verifier_name TEXT,
    required_claims JSONB NOT NULL DEFAULT '[]',
    credential_type TEXT NOT NULL,
    status TEXT NOT NULL DEFAULT 'PENDING'
        CHECK (status IN ('PENDING','COMPLETED','EXPIRED')),
    proof JSONB,
    proof_hash TEXT,
    created_at TIMESTAMPTZ NOT NULL DEFAULT NOW(),
    expires_at TIMESTAMPTZ NOT NULL
);

-- Auto-expire index
CREATE INDEX idx_sessions_expires ON sessions(expires_at);
CREATE INDEX idx_sessions_status ON sessions(status);

-- Verification keys table  
CREATE TABLE verification_keys (
    id UUID PRIMARY KEY DEFAULT gen_random_uuid(),
    credential_type TEXT UNIQUE NOT NULL,
    key_data JSONB NOT NULL,
    created_at TIMESTAMPTZ NOT NULL DEFAULT NOW()
);

-- Row Level Security
ALTER TABLE sessions ENABLE ROW LEVEL SECURITY;
ALTER TABLE verification_keys ENABLE ROW LEVEL SECURITY;

-- Allow relay service role full access
CREATE POLICY "service_role_full_access" ON sessions
    FOR ALL USING (auth.role() = 'service_role');
CREATE POLICY "service_role_keys_access" ON verification_keys
    FOR ALL USING (auth.role() = 'service_role');
\end{lstlisting}

\section{Deployment Configuration}

\subsection{Relay Server (Render)}

The Relay is deployed to Render's cloud hosting platform with the following configuration:

\begin{lstlisting}[language=bash, caption={render.yaml -- Deployment Configuration}]
services:
  - type: web
    name: zeroauth-relay
    env: node
    plan: free
    buildCommand: npm install && npm run build
    startCommand: npm start
    envVars:
      - key: SUPABASE_URL
        fromDatabase:
          name: supabase_url
      - key: SUPABASE_SERVICE_ROLE_KEY
        fromDatabase:
          name: supabase_key
      - key: NODE_ENV
        value: production
      - key: PORT
        value: 3000
\end{lstlisting}

\subsection{Mobile App (Android APK)}

\begin{lstlisting}[language=bash, caption={Android APK Build Command}]
# Development build
eas build --platform android --profile development

# Production APK (local build)
cd android
./gradlew assembleDebug

# APK location
# android/app/build/outputs/apk/debug/app-debug.apk
\end{lstlisting}
